\originalchapter{热平衡与热力学状态}\label{ch:thermo}
\originalsection{系统、状态与温度}
气缸内气体的体积$V$、压强$P$和温度$T$是宏观量; 某个分子的坐标与速度是微观量. 将均匀样品在相同状态下复制一倍, 内能$U$、体积$V$、粒子数$N$和熵$S$加倍, 称为广延量; 温度与压强保持不变, 称为强度量. 表面积、线长和总磁矩也可作为特定系统的广延外参, 但具有边界能或长程相互作用的系统未必严格满足体相可加性.

绝热壁阻止热交换, 导热壁允许热交换. 平衡要求宏观量没有自发净流和变化; 稳态只要求宏观读数不随时间改变, 两端接不同温度热源的导热棒是稳态而非平衡. 固定$N$的简单可压缩物质, 单一均匀相中通常只需两个独立热力学变量. 状态方程如$PV=N\kb T$联系它们, 不能同时把三个量任意指定.

第零定律说, 若A与C热平衡且B与C热平衡, 则A与B也热平衡. 因而可以把状态按热平衡关系分组, 以温度标记这些等温类. 稀薄气体温度计用$PV/N$给经验尺度; 任何严格单调函数都标记相同等温线, 所以第零定律尚未确定开尔文温标. 后文的熵导数或可逆热机效率选定绝对温度, 并以$\kb$联系宏观能量与温度单位.

\originalsection{热、功与第一定律}
采用\textbf{外界对系统做功为正}的约定. 准静态体积功为$\delta W=-P\dd V$: 扩张时系统对外做功, 自身获得的功为负. 拉伸杆的功为$f\dd L$, 增大面积的功为$\gamma_s\dd A$, 电极输送电荷的功为电动势乘电荷增量. 本册用$f$表示张力、$F$表示Helmholtz自由能, 用$\gamma_s$表示表面张力、$S$表示熵, 避免同一个字母承担两个角色. 总磁矩作外参的约定为$B\dd M$; 磁场作外参时需另作Legendre变换, 不能把两套能量混用.

\begin{definition}
内能$U$是系统的状态函数. 热量$Q$是由于温差传递的能量, 功$W$是通过其他宏观操作传递的能量. 第一律为$\dd U=\delta Q+\delta W$. 符号$\delta$提醒我们热与功一般不是状态函数的全微分.
\end{definition}
用绝热过程可测定两个状态的内能差; 同样端点之间的非绝热过程则有不同的$Q$与$W$, 但$Q+W$相同. 闭合循环满足$\oint\dd U=0$, 不意味着$\oint\delta Q$或$\oint\delta W$各为零. 在$P$--$V$图上, $-\int P\dd V$是带符号的功, 路径下的面积直接显示路径依赖.

准静态要求变化足够慢, 使系统可用连续的平衡状态描述. 可逆还要求没有摩擦、有限温差传热等熵产生; 单有“慢”不保证可逆. 气体向真空自由膨胀时系统内压不能作为边界外压使用. 绝热且没有边界功时$\Delta U=0$, 而可逆绝热活塞膨胀会做功, 使内能和温度降低.

\originalsection{全微分与响应函数}
在局部可微状态方程上有倒数关系$(\partial x/\partial y)_z=1/(\partial y/\partial x)_z$及循环关系
\[
\left(\frac{\partial x}{\partial y}\right)_z
\left(\frac{\partial y}{\partial z}\right)_x
\left(\frac{\partial z}{\partial x}\right)_y=-1.
\]
它们来自在$\dd x=x_y\dd y+x_z\dd z$中设$\dd x=0$. 全微分$A\dd x+B\dd y$在单连通光滑区域内, $\partial_yA=\partial_xB$足以保证存在势函数; 不加区域条件时只保证局部存在. 求势时先积分$A$得到$U=\int A\dd x+g(y)$, 再求$\partial_yU$匹配$B$确定$g$. 路径依赖的热或功通常不满足这个相容性条件.

热容必须说明约束: $C_V=(\partial U/\partial T)_{V,N}$, $C_P=(\delta Q/\dd T)_{P,N}=(\partial H/\partial T)_{P,N}$, 其中焓$H=U+PV$. 膨胀率$\alpha=V^{-1}(\partial V/\partial T)_{P,N}$, 等温压缩率$\kappa_T=-V^{-1}(\partial V/\partial P)_{T,N}$, 等熵压缩率$\kappa_S$用$S$代替$T$. 机械稳定的均匀相有$\kappa_T\ge0$. 热膨胀率却可正可负, 水和拉伸橡胶提供实际例子.

先只用第一律, 对准静态体积功有
\[
\delta Q=C_V\dd T+\left[(\partial U/\partial V)_T+P\right]\dd V,
\quad C_P-C_V=\left[(\partial U/\partial V)_T+P\right]\alpha V.
\]
加入第二律后才能消去这里未知的内能导数. 对理想气体, 独立的微观或实验信息给$U=U(T)$, 故$C_P-C_V=N\kb$. 若热容在温区内近似常数, 可逆绝热条件为$C_V\dd T+N\kb T\dd V/V=0$, 积分得到$TV^{\gamma-1}=\text{常数}$和$PV^\gamma=\text{常数}$, $\gamma=C_P/C_V$. 单原子气体$\gamma=5/3$, 含转动振动的气体不能在任意温度沿用这个数值.

声波的压缩周期通常短于局部热扩散时间, 因而近似等熵, 声速满足$c_s^2=(\partial P/\partial\rho)_S=1/(\rho\kappa_S)$. 约束下标在这里有直接实验意义, 不是可以省略的装饰.
